\begin{textsample}{POS Dim 1 – persona\_grok – Score 21.00 – t704\_grok.txt}  \label{ex:f1_pos_033}
World Photography Day, celebrated on August 19, reminds \textbf{us} of the power of \textbf{images} to capture moments that can change the \textbf{world}.
As the legendary \textbf{war} \textbf{photographer} Robert Capa once advised,"If your \textbf{pictures} aren't \textbf{good} enough, you 're not close enough."This wisdom resonates deeply in our \textbf{work} at Greenpeace, where getting close to the issues—literally and figuratively—helps \textbf{us} tell \textbf{stories} that drive environmental action.
I 'm Sudhanshu Malhotra, Greenpeace International 's Multimedia \textbf{Editor} based in Hong Kong.
Photography is profound in its ability to be both deeply personal and publicly \textbf{shared}.
Its origins trace back to the 19th \textbf{century}, when it began \textbf{bearing} \textbf{witness} to the \textbf{world} through landscapes and portraits.
Over the past 200 \textbf{years}, \textbf{images} have been used as evidence in courts and have inspired \textbf{people} to take action.
Photography has long been an ally to environmental \textbf{movements}, communicating across \textbf{languages} and cultures in \textbf{ways} \textbf{words} sometimes cannot.
At Greenpeace, we 've harnessed this power through our \textbf{campaigns}.
Here are 12 iconic \textbf{images} from our archives that highlight the urgency of protecting our planet : - A turtle entangled in plastic waste, illustrating the devastating impact of ocean pollution. - A bloodied \textbf{whale}, an \textbf{image} that helped prompt international whaling bans. - Pelicans soaked in oil, capturing the horror of oil spills. - Vast areas of deforestation in Indonesia, exposing the destruction of vital rainforests. - Protesters at Standing Rock, standing against pipelines threatening indigenous lands and water. - A \textbf{woman} celebrating victory against a coal company, symbolizing community triumphs over fossil fuels.
These \textbf{images}, among \textbf{others} in our collection, show the real faces of environmental crises and the \textbf{hope} in fighting back.
On this World Photography Day, we invite you to join the movement—whether by \textbf{sharing} your own photos, supporting our \textbf{campaigns}, or taking action to protect the planet.
Together, we can create a future worth capturing.

% matched lemmas: bear, campaign, century, editor, good, hope, image, language, movement, other, people, photographer, picture, share, story, us, war, way, whale, witness, woman, word, work, world, year
\end{textsample}
